\documentclass[11pt,a4paper]{article}
\usepackage[utf8]{inputenc}
\usepackage[T1]{fontenc}
\usepackage[margin=2.3cm]{geometry}
\usepackage{graphicx}
\usepackage{booktabs}
\usepackage{xcolor}
\usepackage{hyperref}
\definecolor{copper}{HTML}{C96F4A}
\hypersetup{colorlinks=true, linkcolor=copper, urlcolor=copper}
\setlength{\parskip}{4pt}
\begin{document}
\begin{center}
{\LARGE\bfseries Research Program: The AB Processor (V20, P0–P6)}\\[8pt]
{\large Research Program Monograph · Version 1.0.0}\\[6pt]
{\color{copper}\small RESEARCH PROGRAM · PROCESSOR}\\[4pt]
{\small Isaev Iskhak Khamzatovich · ORCID 0009-0003-7299-0701}\\[2pt]
{\small AB Cloud Quantum Research · DOI 10.5281/zenodo.21825394}
\end{center}

\tableofcontents
\newpage

\section{Research Program: The AB Processor}

\noindent\textit{\color{copper}The processor program turns the mathematics of the core into an executable machine: the AB-cloud ISA, the exact compiler, memory, ALU, zero reading, nonlocality, benchmarks and the passport. Stage P0 is executed; the P0–P6 roadmap is chained with verification.}

\subsection{1. Program goals}

\begin{itemize}
  \item fix the AB-cloud ISA: admissible clouds, parameters, composition;
  \item compile unitary logic exactly: 1−F at the machine epsilon level;
  \item implement "reads, does not carry" algebraically (V23);
  \item provide nonlocality without signaling (V24) and obtain the Datasheet 1.0 passport.
\end{itemize}

\subsection{2. Methodology}

The ISA is not "yet another gate set": every cloud carries a geometric phase, and compilation accounts for it exactly. All stages have pre-registered thresholds and verdict files. Benchmarks follow the E18 methodology: one class, one criterion, strong baselines. numpy replication is mandatory.

\subsection{3. Executed: P0 (V20)}

\begin{table}[htbp]\centering\small
\caption{Table 3.2. The V20 verdict}
\begin{tabular}{p{0.38\textwidth} p{0.58\textwidth}}
\toprule
\textbf{Indicator} & \textbf{Value}
\midrule
accuracy 4×4 & 1−F ≤ 4.4e−16; depth ≤ 7 \\
accuracy 8×8 & 1−F ≤ 6.7e−16; depth ≤ 29 \\
compilation time & 1.8 s for the full run \\
numpy replication & 0 mismatches; |Δ(1−F)| = 2.2e−16 \\
status & PASS; stage P0 closed \\
\bottomrule
\end{tabular}
\end{table}

\subsection{4. The P0–P6 roadmap}

\begin{table}[htbp]\centering\small
\caption{Table 4.1. Processor stages}
\begin{tabular}{p{0.10\textwidth} p{0.13\textwidth} p{0.40\textwidth} p{0.33\textwidth}}
\toprule
\textbf{Stage} & \textbf{Campaign} & \textbf{Content} & \textbf{Depends on}
\midrule
P0 & V20 & ISA + exact compiler & Th5 (block structure) \\
P1 & V21 & volatile memory & ES models, TS1–TS3 \\
P2 & V22 & AB-cloud ALU & V20 \\
P3 & V23 & zero reading & E17, Th6 \\
P4 & V24 & nonlocality & E13/E14 \\
P5 & V25 & benchmarks & the E18 methodology \\
P6 & V26 & Datasheet 1.0 & all verdicts of the P-chain \\
\bottomrule
\end{tabular}
\end{table}

\subsection{5. Honest boundaries}

The 2q KAK optimum (3 gates) is unreachable in the ISA class: the compiler produces 10, and this is a property of the class, not of the algorithm. The V20 targets are unitary 4×4/8×8; non-unitary operators and larger sizes are outside the campaign. The V23/V24 previews are not full verdicts: the full campaigns are ahead.

\subsection{6. Key results in detail}

\subsection{6.1. From P0 to the P3/P4 previews}

\subsubsection{V20 — exact compilation (P0)}

The only executed processor campaign so far closes stage P0: the ISA is fixed, the Givens-QR compiler assembles unitary 4×4 and 8×8 from AB clouds with infidelity 1−F ≤ 4.4e−16 and 6.7e−16, depths ≤ 7 and ≤ 29 gates, the full run takes 1.8 s. The numpy replication gives zero counter mismatches and |Δ(1−F)| = 2.2e−16. Accuracy at machine epsilon means: further improvement requires changing the arithmetic, not the algorithm.

\subsubsection{V23 preview — zero reading}

The P3 stage preview: the processor receives six constants of the trace formula and reproduces 12 of 12 ζ zeros exactly; storage = 0. The blind discipline is preserved: the constants contain no zeros, the output is pinned before validation. A preview is not a full verdict: the full V23 campaign requires the whole E17 methodology at the processor level.

\subsubsection{V24 preview — nonlocality}

The P4 stage preview: a pair of separated clouds gives CHSH S = 2√2 − 4.4e−16, signaling protection 1000/1000, zero commutators of the measuring operators. The correlations are nonlocal but carry no signal — a structural property of quantum mechanics reproduced by the architecture and required for distributed operations.

\subsubsection{E18 link — class B}

The external control of the processor compiler: against five baselines at equal ISA — a ×59,000 time speed-up at equal accuracy; the honest boundary — the 2q KAK optimum (3 gates) is unreachable in the class (the compiler produces 10). Class B protects the comparison from looseness and makes the result reproducible.

\subsection{7. How to reproduce}

\begin{verbatim}
python3 modules/processor_ab/scripts/compute_v20.py      # campaign V20
cd modules/processor_ab && python3 -m pytest tests/ -q   # tests
# web lab: modules/processor_ab/web_lab/index.html (compiler tab)
\end{verbatim}

\subsection{8. Verdicts and data}

\begin{itemize}
  \item V20 verdict: modules/processor\_ab/data/v20\_verdict.json;
  \item claim audit: modules/processor\_ab/verification/claim\_audit/;
  \item figures: figures\_hires/ru|en (600 dpi); workbook: tables/qp\_workbook.xlsx;
  \item rpbridge: ISA, compiler, ALU, reader — modules/processor\_ab/rpbridge/.
\end{itemize}

\subsection{9. Program risks}

The main risk is extending the ISA without re-examining the class: any new primitives must pass the E18 methodology again. The second risk is the temptation to claim the KAK optimum: the class boundary is fixed by a verdict and is not removed without new verdicts. The third — the V23/V24 previews: until the full campaigns, their results do not enter the passport.

\end{document}
